\documentclass[10pt]{article}
\usepackage[margin=2.6cm]{geometry}
\usepackage{amsmath,amssymb,amsthm}
\usepackage{booktabs}
\usepackage[hidelinks]{hyperref}
\title{\bfseries The plateau versus $2\pi e$}
\author{Denis Joubert\\ \small Independent researcher, Switzerland
\and Grok\\ \small draft for the author, 1 September 2026}
\date{September 1, 2026 --- $\Delta_\infty$ lock}
\begin{document}
\maketitle
\begin{abstract}
\noindent
The published plateau is $\Delta_\infty\approx16.0\pm1.6$ on
$\ell\in[40,80)$ and $16.0\pm2.1$ on $[80,150)$. The constant
$2\pi e\approx17.079$ sits inside both bands. New $\chi_3$ ladders
at $\mu=16,22,30$ climb $15.88\to16.92\to17.07$ at levels
$40,62,105$. The deepest new spacing matches $2\pi e$ to $0.01$ nat.
The lock ``the plateau stays at $16$'' is dead. $2\pi e$ survives;
one character and one basis do not select it.
\end{abstract}

\section{The lock}

A1 said the leftover Rayleigh problem drops by a universal factor
once the desert is spent. The published factor is $e^{-16}$. Among
the small list of constants that could have been that factor,
$2\pi e\approx17.079$ is the only one inside the plateau bands
and not already killed ($s=\gamma_1^2/(2\pi e)$ died on $\chi_4$).

\section{New rungs}

\texttt{code/scan\_s.py} on $\chi_3$ (odd, $s=4$, depth-adequate
at these windows). Spacings $\ell_k-\ell_{k+1}$ at the deep end:

\begin{center}\small
\begin{tabular}{lccc}
\toprule
$\mu$ & $N$ & deepest $(\ell,\Delta)$ & $2\pi e-\Delta$ \\
\midrule
16 & 37 & $(39.5,\,15.88)$ & $+1.20$ \\
22 & 41 & $(62.3,\,16.92)$ & $+0.16$ \\
30 & 45 & $(104.7,\,17.07)$ & $+0.01$ \\
\bottomrule
\end{tabular}
\end{center}
The sequence climbs toward $2\pi e$, not toward $16$. At $\mu=30$
the match is $17.07$ against $17.079$.

The rest of each ladder still follows the published ramp
($9$--$13$ nats on $\ell<40$). Only the deepest spacing is the
plateau test.

\section{Second character, next window}

\texttt{scan\_s.py} omitted the primes $37+$ --- the accidental
quorum of notebook 16.4. With the sieve extended and eigenvalues
sorted, at $N=53$:

\begin{center}\small
\begin{tabular}{llcc}
\toprule
$\chi$ & $\mu$ & deepest $(\ell,\Delta)$ & $2\pi e-\Delta$ \\
\midrule
$\chi_3$ & 38 & $(128.8,\,16.82)$ & $+0.26$ \\
$\chi_4$ & 30 & $(81.4,\,17.65)$ & $-0.57$ \\
$\chi_4$ & 38 & $(103.7,\,18.16)$ & $-1.08$ \\
\bottomrule
\end{tabular}
\end{center}
$\chi_3$ hovers $16.8$--$17.1$ across four windows. $\chi_4$ sits
$0.6$--$1.1$ nats above $2\pi e$. The plateau is the interval
$[16.8,18.2]$, not a selected constant. $2\pi e$ remains inside
the interval; it is no longer the unique occupant. The $N=75$
lock is still the next stroke, not a requirement to kill $16$.
\end{document}
